\section{From the proofs to Lean derivations}\label{sec:formal}
The formalization target for a closed arithmetical claim $\varphi$ is a value of type $[\,]\vdash[\mathsf{PRAAx}]\,\varphi$ in ProveIt's ordinary system. For a schema, first fix the PR syntax of the evaluators; put their arithmetic contracts in the antecedent or context. A semantic Lean theorem about functions on \texttt{Nat} is a useful specification. To certify the strength of the theorem one must construct the indicated object-language derivation.

\subsection{Syntax, modules, and a worked target}
We used the following interfaces from the repository at commit
\href{https://github.com/VladimirReshetnikov/ProveIt/tree/6d03f3643ebb91a6243eb1e35e7599b388197786}{\texttt{6d03f3643e}}.
\texttt{ConciseIPRA.PR} defines \texttt{PR k} (zero, successor, projections, composition and primitive recursion), terms \texttt{PTm n} built by variable and application, and \texttt{PRAAx}. In \texttt{PR.prec} the recursion argument occurs first; the step arguments begin with the predecessor and old state. The formulas and derivations use \texttt{Fm} and \texttt{Prf}. \texttt{PRA.qfInd} takes the quantifier-free check and derivations for the base and step. The module \texttt{IPRADerivations.Foundations.GuardedTrace} supplies \texttt{guarded\_loop}, \texttt{guarded\_trace}, and versions extracting every stage. \texttt{IPRADerivations.Core.IntCertificates} supplies signed integer terms \texttt{ZT}, their equality formula, and the sound ring lifting operation \texttt{ZT.derive}. Use \texttt{pra\_iff\_native} in \texttt{ConciseIPRA.NativeBiInterpretation} once an ordinary derivation has been obtained.

\paragraph{The rectangular target from the circle argument.}
Encode horizontal rational increments on a $(w,h)$ grid by a row-major array $H$ with $w(h+1)$ entries. Encode the vertical ones by $V$ with $(w+1)h$ entries. Multiply by the product of the positive denominators, so both signed arrays use integers and a common positive scale $D$. Index $H$ by $(j,i)\mapsto jw+i$, $V$ by $(j,i)\mapsto j(w+1)+i$. Let $B(j,r)$ be the sum of the first $r$ horizontal entries of row $j$. The displayed formula is an actual target for numerical coding, with each bounded conjunction expanded as in Lemma~\ref{lem:bounded}:
\[
\begin{split}
 \operatorname{Valid}(H,V,w,h)\ \land\;
 &\bigwedge_{j<h,\ i<w}
 [H_{ji}+V_{j,i+1}=V_{ji}+H_{j+1,i}]\\[-2pt]
 {}\land\;&\bigwedge_{j<h}[V_{jw}=V_{j0}]
       \quad\Longrightarrow\quad B(0,w)=B(h,w).
\end{split}
\]
Here subscripts such as $V_{jw}$ mean row $j$, column $w$. The inner invariant at prefix $r$ is
\[
 B(j,r)+V_{j,r}=V_{j,0}+B(j+1,r).
\]
At zero it is an identity. Add the next cell equality and cancel the repeated vertical edge. At $r=w$ substitute seam equality. The outer invariant $B(0,w)=B(j,w)$ advances by transitivity. Both are quantifier-free equations after substituting the PR table-access and summation terms. This is the entire uniform telescope derivation, including the empty rectangle. To finish the circle theorem, add a stage selecting signed residues in $(-D/4,D/4)$ at each edge; its bounded flag records congruence to the difference of endpoint readings modulo $D$. A cell sum has absolute value $<D$ and is divisible by $D$, hence zero. On the bottom row take the common scale divisible also by $N$ so the total is $D$. On the constant top row it is zero. The side and sampling assumptions in Lemma~\ref{lem:winding} supply exactly those checks.

The new file \texttt{lean/GridWinding.lean} expresses both iterations for ordinary Lean integer arrays. \texttt{lean/ObjectGrid.lean} gives a candidate proof-producing version: for every external rectangle size and every array of symbolic \texttt{ZT} entries it combines object-level cell proofs and seam proofs into an object-level row equality. The last step for the \emph{uniform} numerical target is to replace the two external size inductions by the two object-level guarded loops just specified. These new files were not compiled in this build. The repository's reusable finite-topology, metric, algebraic and sphere modules are included unchanged in \texttt{companions/}; their upstream build status is recorded separately in the archive guide.

\subsection{Definitions and order of the remaining work}
The relevant definitions, bounds, and witness-producing invariant for each family occur together in the preceding proofs. The following is a dependency order for implementing them in the same object calculus. All functions can be made total by returning zero on malformed tables; validity always occurs in the antecedent.
\begin{enumerate}
\item \textbf{Arrays and arithmetic.} Define \texttt{bit}, \texttt{read}, \texttt{write}, \texttt{firstHit}, \texttt{allBelow}, prefix parity, and signed rational comparisons using the encoding in Section~\ref{sec:arithmetic}. Prove read-after-write at equal and unequal indices, the prefix fold equations, and the witness specifications of Lemma~\ref{lem:bounded}. For rational linear operations prove compatibility with scaling both numerators and denominators and enforce the denominator tests.
\item \textbf{Spaces and graphs.} \texttt{Open} tests upward closure over all $n^2$ pairs. \texttt{LeastNbr} intersects all open masks containing the given point. \texttt{Reach} repeats one mark expansion $n$ times, storing a parent at each new marked vertex. Prove that marks have recorded paths and that closure under adjacency holds at the end. The component, subcover, quotient and beat-point verifiers in the finite section use these three operations. For grid drawings define the two edge arrays, row parity and the dual adjacency test. The exact recurrences of Lemma~\ref{lem:gridfill} yield boundary after filling, and filling after boundary; the forest proof gives the face count.
\item \textbf{Chains and reduced words.} Define face size, simplex image and boundary by enumeration of one-vertex deletions. Pair the two deletion orders; each target bit then cancels in $d^2$. Carry a column and its preimage together in elimination. Extend this in the order specified at the end of Section~\ref{sec:algtop}; the projection matrices from interval splitting give the barcode. For the polygon models use the explicit signed substitutions, testing the free coordinates and the residue modulo two. For graph words use a signed-letter stack. Store the current representative and emitted symbols for Theorem~\ref{thm:schreier}. Its prefix invariant is that their expansion equals the scanned input prefix times the inverse of the representative. In folding, store the vertex equivalence and images of oriented edges; the two-edge bridges lift a given word with a bounded certificate at each fold. For Fox colorings the data and move bijections are the crossing records of Section~\ref{sec:knotcolor}.
\item \textbf{Cauchy names and charts.} For each fixed input array term put the universally quantified regularity and modulus formulas in the context. A new name is an actual PR term in the precision and all free parameters. To prove a comparison, fix all outer stage indices, derive a rational estimate with auxiliary error $C2^{-s}$ and remove that error as in the metric section. Implement grid rounding, clamped interval maps, rational square-root approximation, sign-certified division, and the local chart inverses there. Cantor coding adds finite binary prefixes, ternary interval endpoints and finite metric nets. Branch decisions on real names are implemented through finite readings with a proved positive margin, or by gluing uniformly compatible rational formulas.
\item \textbf{Higher homotopy and surfaces.} Implement \texttt{kuhnWeights}, deterministic rational elimination for a covector, root counting by polynomial division and the $j$-scan on the moment curve. The output certificate contains the successful $j$, all covectors and the rational lower bound $b$. Check annihilation and nonzero pairings by finite products before defining the normalized coning name. The collar formulas prove the boundary equations at every stage. For the Möbius atlas, store $m$ seam bits; the invariant for a lifted side after $i$ seams is $s\prod_{j<i}e_j$. The branch obstruction uses just three segment maps and the interval search, with no classification of general manifolds.
\end{enumerate}
\paragraph{What is checked.}
The archive provides runnable rational and finite regression checks for the principal combinatorial constructions. The supplied Lean code and declaration-level plan give a path to formal completion. The complete collection of uniform object-language derivations is still to be constructed. The bare theorem that a search routine is PR, a passing finite test, or an ordinary Lean proof in a stronger logic does not close that remaining step.
